\chapter{Conclusion and Future Work}

\section{Conclusion}

This project set out to investigate, implement, and evaluate dynamic partial reconfiguration (DPR) of a RISC-V co-processor on a commodity FPGA development board, using open-source and vendor toolchains as the primary implementation vehicles. The objectives spanned three distinct layers of inquiry: a literature-level analysis of state-of-the-art DPR toolflows, a hands-on baseline verification of the AMD Vivado Dynamic Function eXchange (DFX) methodology, and a functional proof-of-concept demonstrating runtime swap of a custom function unit (CFU) attached to a VexRiscv RISC-V soft-core processor.

\subsection{Literature and Toolflow Analysis}

Three major toolflows were surveyed at depth: RapidStream, Hierarchical Partial Reconfiguration (HPR), and ZyPR/ZyCAP. Each addresses a different dimension of the DPR problem space. RapidStream targets compilation latency by decomposing a design into pipeline-latency-insensitive floorplan islands and compiling them in parallel, achieving order-of-magnitude reductions in compile time relative to monolithic Vivado/Vitis flows. HPR extends this philosophy to the runtime domain, constructing recomposable reconfigurable pages connected through a Hoplite-BF-variant Network-on-Chip, thereby removing the need for manual operator-to-floorplan assignment. ZyPR and its associated ZyCAP runtime controller represent the most complete end-to-end realisation surveyed: ZyCAP achieves partial bitstream throughput of up to 400 MB/s on Zynq platforms through AXI4 DMA into the configuration port, with LRU caching of previously loaded modules in DRAM.

Hands-on reproduction of all three toolflows was attempted. RapidStream could not be evaluated: the toolflow has been transitioned into a proprietary commercial product, and all public API keys for the academic release have been revoked. HPR reproduction was attempted but ultimately abandoned due to version incompatibilities between the required Vivado revision, the available Linux kernel, and associated library dependencies — combined with the compute constraints of the development machine, which was unable to sustain the place-and-route workload without job failure. ZyPR similarly could not be exercised, as the toolflow depends on a proprietary component requiring an educational licence that was not granted. These three toolflows are therefore covered in this report through rigorous literature analysis rather than experimental reproduction. Their architectural concepts, however, directly informed all design decisions made during the implementation phases.

\subsection{DFX Baseline Validation}

The AMD UG947 DFX tutorial was successfully implemented on the Artix-7 xc7a100tcsg324-1 device (Arty A7-100T) as a baseline validation of the Vivado DFX toolflow. Two interchangeable configuration pairs — \texttt{shift\_left/count\_down} and \texttt{shift\_right/count\_up} — were synthesised, placed, routed, and verified to produce correct partial bitstreams. Timing closure was achieved in both configurations, with a Worst Negative Slack of +2.022 ns against a 200 MHz clock constraint, and with all hold endpoints satisfied at a WHS of +0.109 ns. Resource utilisation was minimal by design (10 LUTs, 64 flip-flops, 1 BRAM tile), confirming that the toolflow infrastructure — pblock constraints, configuration run hierarchy, and partial bitstream generation — is correctly set up and produces bitstreams that can be loaded onto hardware without error. This exercise established the essential foundation of familiarity with Vivado's DFX design flow that was subsequently carried into the CFU Playground integration.

\subsection{Reconfigurable RISC-V CFU System}

The primary contribution of this project is the implementation and functional demonstration of a dynamically reconfigurable Custom Function Unit (CFU) attached to a VexRiscv RV32I soft-core processor embedded within a LiteX SoC, targeting the Arty A7-35T (xc7a35ticsg324-1L). The full system was synthesised and implemented using Vivado 2024.2, consuming 4,311 LUTs (20.73\%), 6,608 flip-flops (15.88\%), and 5.5 BRAM tiles (11\%) — leaving substantial fabric headroom for future expansion. The design instantiates the ICAPE2 configuration primitive directly in the programmable logic fabric, occupying one of the device's two available ICAP sites, and monitors the End-of-Startup (EOS) signal through the STARTUPE2 primitive to provide hardware-timed measurement of the reconfiguration window via the \texttt{recon\_counter} module.

Two CFU configurations were developed and verified: one implementing integer multiplication and one implementing bitwise manipulation operations. A critical behavioural property central to the safety of any DFX co-processor system was confirmed experimentally: when an instruction carrying a \texttt{function\_id} that does not match the currently loaded CFU configuration is issued by the VexRiscv pipeline, the DFX Decoupler holds the static-side CFU response bus at its default values, producing a zeroed output. The processor receives this result and continues execution without stall or exception, demonstrating that the decoupler correctly isolates the processor from an incompatible or transitioning reconfigurable module.

Reconfiguration latency was measured across two platforms and three interfaces. On the Arty A7-35T, JTAG-based reconfiguration of the 164,534-byte partial bitstream was measured at 63.19 ms (6,318,898 clock cycles at 100 MHz, captured by the hardware \texttt{recon\_counter}). On the PYNQ-Z2, JTAG latency for the 203,308-byte partial bitstream was measured at 62.8 ms, while PS-side PCAP reconfiguration achieved 1.1 ms at a throughput of 183.2 MB/s — a 57× reduction in reconfiguration latency compared to the JTAG path. Across both boards, JTAG-based partial reconfiguration was successfully demonstrated and validated, confirming the fundamental DPR mechanism.

The project thus demonstrates the feasibility of runtime RISC-V CFU reconfiguration on commodity 7-series FPGA hardware using open toolchains, establishes a quantitative reconfiguration latency baseline, and validates the DFX Decoupler as a reliable isolation mechanism between the static processor and the reconfigurable compute fabric.

\vspace{0.4em}\hrule\vspace{0.4em}

\section{Future Work}

The current implementation represents a functional first iteration that validates the core DPR concept. Several well-defined development paths remain open, each of which addresses a specific limitation of the present system.

\subsection{Autonomous Partial Bitstream Delivery via SPI Flash}

The most immediate and highest-priority extension is replacing the host-driven JTAG reconfiguration path with an autonomous on-fabric bitstream fetch. The Arty A7-35T carries a 128 Mbit (16 MB) Quad-SPI flash device, which is already accessible through the \texttt{axi\_quad\_spi\_0} peripheral instantiated in the current design (it contributes 6 mW to the power budget in the current build, confirming it is live). The intended flow is for the VexRiscv firmware to issue a software command — via a custom CSR or a memory-mapped register write to the DFX controller — that causes the static fabric to read a partial bitstream from a pre-determined flash address and stream it word-by-word into the ICAPE2 primitive, without any host intervention.

This path is architecturally straightforward: the AXI Quad SPI controller, the DFX controller IP, and the ICAPE2 are already instantiated and connected. The remaining work is completing the state machine in the \texttt{system\_wrapper} that sequences the SPI read, the ICAP write burst, and the decouple/reconnect handshaking with the DFX Decoupler. Once this is functional, the system will be capable of fully autonomous CFU swapping from firmware, closing the loop between the RISC-V instruction stream and the reconfigurable fabric.

\subsection{DDR-Backed Runtime Bitstream Buffering}

Once SPI flash fetch is operational, the logical next step is to investigate DDR-based bitstream storage for lower-latency random access. Two approaches are under consideration. The first is a boot-time migration strategy: at power-on, the firmware copies all available partial bitstreams from flash into DDR, after which runtime selection becomes a DDR read rather than a slower SPI access. The second approach transfers bitstreams directly from the host to DDR at runtime via a UART-to-AXI DMA path, enabling dynamic loading of new CFU configurations without reflashing the device.

This line of work also represents the path towards a ZyCAP-style architecture on pure PL fabric: a high-throughput DMA engine reading from DDR and writing to the ICAPE2 at rates approaching the 400 MB/s theoretical ceiling of the 7-series ICAP interface. At that throughput, the 164,534-byte Arty partial bitstream would transfer in under 0.5 ms, making per-task CFU swapping viable even in moderately time-constrained embedded applications.

\subsection{ICAP Throughput Characterisation and Comparison with PCAP}

The current \texttt{recon\_counter} module provides cycle-accurate hardware timing of the JTAG reconfiguration window. Once the autonomous ICAP path is implemented, the same counter will measure the ICAP reconfiguration window, producing an experimentally grounded ICAP throughput figure for the Arty A7-35T. This measurement can then be directly compared against the PCAP throughput measured on the PYNQ-Z2 (183.2 MB/s), providing a meaningful architectural comparison between the PS-assisted Zynq approach and the pure-PL Artix-7 approach. Resolving the ICAPE2 timing constraint mismatch identified in Chapter 6 — by applying a multicycle path exception or by clocking the ICAP interface from a dedicated divided domain — would also be a prerequisite for a clean, production-quality timing report.

\subsection{Multiple Simultaneous Reconfigurable Partition Slots}

The current architecture defines a single reconfigurable partition (\texttt{cfu\_compute}) that holds one CFU configuration at a time. A natural architectural extension is to define two or more independent pblocks, each of which can host a different CFU module simultaneously. This mirrors the HPR concept of multiple recomposable PR pages and the ZyPR O1 optimisation level, in which multiple operators are mapped to fixed PR pages and can be independently swapped at runtime without affecting one another. With approximately 75–80\% of the Arty A7-35T fabric remaining unused after the current implementation, the headroom exists to support at least two additional partition slots of comparable size to \texttt{cfu\_compute}. A software scheduler running on the VexRiscv could then assign incoming workloads to the partition whose loaded configuration matches the required function, falling back to software emulation for mismatched operations while a background reconfiguration proceeds.

\subsection{Graph-Based CFU Selection and Scheduling}

As the number of available CFU configurations grows, static assignment of bitstreams to tasks becomes impractical. A graph-based runtime scheduler — analogous to the task-graph approach explored in the DML work surveyed in Chapter 2 — could be implemented in the VexRiscv firmware to track which configuration is currently loaded in each partition, predict the reconfiguration cost of a proposed swap, and make quantitative decisions about whether to swap, wait, or fall back to software. This would transform the current proof-of-concept into a genuinely policy-driven reconfigurable compute system.

\subsection{Extension to RV32IMC and More Complex Accelerators}

The current VexRiscv instance is configured as a minimal RV32I core. Enabling the M extension (integer multiplication and division instructions) and the C extension (compressed instructions) within VexRiscv and re-synthesising the SoC would increase the baseline instruction throughput available to the firmware. Against this extended baseline, more sophisticated CFU accelerators — fixed-point DSP filters, small neural-network inference engines, or cryptographic co-processors — could then be implemented as reconfigurable modules and their contribution measured in terms of cycles saved per operation relative to the M-extension software path. This would complete the performance validation dimension of the project that the current proof-of-concept, which focuses on correctness of mechanism rather than quantified speedup, intentionally defers.

\subsection{Zynq ICAP Integration}

On the PYNQ-Z2, the Zynq-7020's PL fabric contains an ICAP primitive that is physically accessible but requires careful coordination with the PS to avoid conflicts with the PCAP interface, as both share the configuration port arbitration logic. Implementing a PL-side ICAP controller on the PYNQ-Z2 — using the same \texttt{system\_wrapper} and ICAPE2 instantiation approach developed for the Arty A7-35T — would enable a direct comparison of ICAP throughput versus PCAP throughput on the same device, removing the platform variable from the comparison. It would also demonstrate that the pure-PL DPR architecture is portable across Artix-7 and Zynq-7000 devices without architecture-specific modification.

\vspace{0.4em}\hrule\vspace{0.4em}

\section{Summary}

This project has demonstrated that runtime reconfiguration of a RISC-V custom function unit is achievable on low-cost 7-series FPGA hardware within the constraints of the AMD Vivado DFX toolflow and the CFU Playground framework. The DFX Decoupler was confirmed to provide reliable isolation between the static RISC-V processor and the reconfigurable compute partition, and the \texttt{recon\_counter} hardware module provided cycle-accurate timing that validated the JTAG reconfiguration latency measurement. PCAP-based reconfiguration on the Zynq platform demonstrated the 57× latency advantage of high-bandwidth programmatic configuration over JTAG, establishing a clear performance target for the ICAP-based autonomous flow that constitutes the immediate next development phase.

The gap between the current state — functional JTAG-triggered reconfiguration with manual bitstream delivery — and the target state — fully autonomous, firmware-triggered CFU swapping from on-board non-volatile storage — is well-defined and bounded. The architectural primitives (ICAPE2, AXI Quad SPI, DFX controller, DFX Decoupler) are all present and functionally instantiated in the current bitstream. The remaining work is principally a firmware and RTL integration exercise rather than an architectural redesign, which positions the system for rapid completion of the autonomous reconfiguration loop in the next iteration of this work.
